\section{The actual Haar density in exponential coordinates}
\label{sec:haar-chart}
The chart factor in the packet is the density of probability Haar
measure. We construct it from the given local exponential and logarithm,
keeping its scalar normalization. Write $\mathfrak g$ for the physical
colour space, with its counting Lebesgue measure, and $\exp$ for the
group exponential. The local logarithm is inverse to $\exp$ near zero.
After shrinking the radius, $\exp$ is a homeomorphism from a ball
$D$ about zero onto an open neighbourhood $U$ of the identity. To
check openness, if $b$ is in a sufficiently small open set $O$, continuity
of the local logarithm gives a neighbourhood of $\exp b$ whose
logarithms lie in $O$; the local right-inverse identity places that
neighbourhood in $\exp O$.

For a group element $u$, where the expression stays in the chart set
\[
 T_u(b)=\log(u\exp b),\qquad
 H(b)=D T_{(\exp b)^{-1}}(b),\qquad J(b)=\det H(b).
\]
The map $(b,c)\mapsto\log((\exp b)^{-1}\exp c)$ is analytic
near every point $(b,b)$: its matrix argument is the identity there,
where the specified logarithm is analytic. Differentiating in $c$
and restricting to the diagonal shows that $H$, and hence $J$, is
analytic. At zero, $T_1$ agrees with the identity on a neighbourhood,
so $H(0)=1$ and $J(0)=1$. Shrink $D$ so that
$1/2<J(b)<2$ throughout $D$.

\subsection{The overlap identity and normalization}
If $b,T_u(b)\in D$, put $y=T_u(b)$. Locally at $b$,
\[
 T_{(\exp b)^{-1}}
       =T_{(\exp y)^{-1}}\circ T_u.
\]
Indeed $\exp y=u\exp b$, and the intermediate exponential and
logarithm cancel on this overlap. Differentiate at $b$ and take
determinants to obtain
\begin{equation}\label{eq:haar-overlap}
 J(T_u(b))\det D T_u(b)=J(b).
\end{equation}
Both densities are positive, so the transition determinant is positive
as well.

Let $\nu=\exp_*(1_DJ(b)\,db)$. It is a nonzero finite measure,
supported on $U$: positivity gives nonzero mass, and the upper bound
$J<2$ on the bounded ball gives finiteness. For a measurable
$A\subset U$, change variables by $T_u$ on
$\{b\in D:u\exp b\in A\}$. This map is injective there and has
image $\{y\in D:\exp y\in A,\ u^{-1}\exp y\in U\}$;
its inverse is $T_{u^{-1}}$. Equation~\eqref{eq:haar-overlap}
therefore gives
\begin{equation}\label{eq:haar-local-translation}
 \nu(u^{-1}A)=\nu(A\cap uU).
\end{equation}
All sets lie in open chart domains, so the ordinary finite-dimensional
change-of-variables formula applies also to their measurable subsets.

Let $h$ be probability Haar measure on the compact group. Integrate
\eqref{eq:haar-local-translation} over $u$ with respect to $h$.
Tonelli's theorem and right invariance give
\[
 \int\nu(u^{-1}A)\,dh(u)
 =\int\!\int 1_A(ux)\,dh(u)\,d\nu(x)
 =\nu(G)h(A).
\]
On the other side, inversion and right translation preserve $h$, so
\[
 \int\nu(A\cap uU)\,dh(u)
 =\int 1_A(x)\!\int1_U(u^{-1}x)\,dh(u)\,d\nu(x)
 =h(U)\nu(A).
\]
The open nonempty set $U$ has positive Haar measure. Define
$c_{\rm chart}=h(U)/\nu(G)>0$. Then, for every measurable $A\subset U$,
$h(A)=c_{\rm chart}\nu(A)$. Since $\nu$ is supported on $U$, this
is the exact identity of measures
\begin{equation}\label{eq:haar-chart-measure}
 \exp_*\bigl(1_D c_{\rm chart}J(b)\,db\bigr)=h|_U.
\end{equation}
This argument constructs both the analytic density and its scalar
normalization from the local chart and the given probability Haar
measure. It makes no assertion that additive Lebesgue measure is a
Haar measure on the group.

\subsection{The protected chart correction}
Fix $a>0$ and use the prescribed small scale
$t=P(h)/\sqrt h>0$ from \eqref{eq:warning-scales}. It has upper
exponent $-1/2$, and $t^{-1}\le\sqrt h$ because $P(h)\ge1$;
thus its reciprocal has upper exponent $1/2$. Let $\chi:[0,\infty)\to[0,1]$
be the smooth cutoff, equal to one on $[0,1]$ and zero on $[4,\infty)$.
For a fluctuation bond $u_i$ set
\[
 \gamma_a(u_i)=\chi\!\left(\frac{\|u_i\|^2}{a^2t^2}\right),
 \qquad j_i(u)=\gamma_{2a}(u_i)\log J(u_i).
\]
Outside the chart this last expression can use any specified measurable
extension of $J$ and the total real logarithm convention; the cutoff
makes its value zero there once $4at$ is smaller than the chart radius.
Where $\gamma_a(u_i)\ne0$, we have $\|u_i\|\le2at$, and hence
$\gamma_{2a}(u_i)=1$. Positivity of $J$ then proves
\begin{equation}\label{eq:haar-chart-exact}
 e^{j_i(u)}=J(u_i)\quad\hbox{where }\gamma_a(u_i)\ne0.
\end{equation}
Thus the doubled cutoff preserves the true Haar factor on the support
of the narrower protected integral.

Since $\log J$ is analytic near zero and vanishes there, the
fundamental theorem of calculus on a smaller closed ball gives
$|\log J(v)|\le C\|v\|$. If $\gamma_{2a}(u_i)$ is nonzero,
$\|u_i\|\le4at$; if it is zero, $j_i=0$. Therefore, globally in $u$,
\begin{equation}\label{eq:haar-chart-value}
 |j_i(u)|\le4aCt.
\end{equation}
In particular the bound $|j_i|\le1$ used in the packet estimate holds
at one sufficiently large coupling threshold, uniformly in the period
and in the number of chart bonds.

We record the differentiated version with its exact incoming premise.
Suppose a variation $u_i(y)$ is $C^k$ near $y_0$, and its increment
$u_i(y)-u_i(y_0)$ has derivatives through order $k$ bounded by
$v(x)$, where $v$ has upper exponent $-1/2$. This is the compensated
variation supplied by the Gaussian change of variables; its production
is separate from the chart calculation here. If $\gamma_b(u_i(y))$
is zero on a neighbourhood of $y_0$, all its jets vanish. Otherwise
continuity gives $\|u_i(y_0)\|\le2bt$. Thus every derivative,
including order zero, of $u_i$ at $y_0$ is bounded by $v+2bt$.
For $g=u_i/(bt)$ these bounds are at most
\[
 K_b=(bt)^{-1}(v+2bt),
\]
which has upper exponent zero. The derivatives of $\|g\|^2$ are
bounded by a fixed polynomial in $K_b$; this follows directly by the
product rule for the inner product. Derivatives of $\chi$ are bounded
on $[0,4]$, and its value is at most one. The chain rule, indexed by
partitions of the derivative arguments, therefore gives
\begin{equation}\label{eq:haar-cutoff-jets}
 \|D^r\gamma_b(u_i(y_0))\|\le L_b(x)^r,
 \qquad 0\le r\le k,
\end{equation}
for an $L_b\ge1$ of upper exponent zero. The zero-order inequality
uses the range of $\chi$; a uniform amplitude from the positive-order
bounds can be absorbed into $L_b$.

Apply this to the broader cutoff $b=2a$. If its germ vanishes, so do
all jets of $j_i$. Otherwise $\|u_i(y_0)\|\le4at$, and all jets
of $u_i$ are bounded by $V=v+4at$, which tends to zero. All derivatives
of $\log J$ through order $k$ are bounded on a fixed smaller chart
ball. For positive derivative order, each partition term in the
chain rule contains at least one derivative of $u_i$; since $V\le1$,
its product is at most $V$. At order zero use
$|\log J(u_i(y_0))|\le CV$. Hence
\[
 \|D^r(\log J\circ u_i)(y_0)\|\le C_kV
 \quad(0\le r\le k).
\]
Leibniz's formula with \eqref{eq:haar-cutoff-jets} yields
\begin{equation}\label{eq:haar-chart-jets}
 \|D^r j_i(y_0)\|\le C_kV(1+L_{2a})^k
 \quad(0\le r\le k).
\end{equation}
The common right-hand side has upper exponent $-1/2$. The argument
covers points on the mask boundary and arbitrary Gaussian inputs:
smallness is inferred from the nonvanishing germ of the protected
factor, not assumed for the Gaussian field globally.
